\documentclass[11pt,a4paper]{article}

% ============================================================
% REEM V2 — INVENTORY INTELLIGENCE REPORT
% ============================================================

\usepackage{fontspec}
\usepackage{polyglossia}
\setmainlanguage{english}
\setmainfont{TeX Gyre Pagella}

\usepackage{geometry}
\geometry{margin=1in}

\usepackage{amsmath}
\usepackage{amssymb}
\usepackage{booktabs}
\usepackage{longtable}
\usepackage{tabularx}
\usepackage{array}
\usepackage{graphicx}
\usepackage{float}

\usepackage{tikz}
\usetikzlibrary{arrows.meta, positioning, shapes}

\usepackage{sectsty}

\usepackage{pgfplots}
\usepackage{pgfplotstable}
\pgfplotsset{compat=1.18}

\usepackage{xcolor}
\usepackage{microtype}
\usepackage{enumitem}
\usepackage{hyperref}
\usepackage{fancyhdr}
\usepackage{titlesec}
\usepackage{tcolorbox}

% ============================================================
% COLOURS
% ============================================================

\definecolor{titleColor}{HTML}{800000}
\definecolor{blueColor}{HTML}{4682B4}
\definecolor{purpleColor}{HTML}{6A5ACD}
\definecolor{goldColor}{HTML}{DAA520}

\hypersetup{
    colorlinks=true,
    linkcolor=blueColor,
    urlcolor=blueColor
}

\sectionfont{\color{titleColor}}
\subsectionfont{\color{blueColor}}

\pagestyle{fancy}
\fancyhf{}
\fancyhead[L]{REEM V2}
\fancyhead[R]{Inventory Intelligence}
\fancyfoot[C]{\thepage}

% ============================================================
% DOCUMENT
% ============================================================

\begin{document}

\begin{titlepage}

\centering

\vspace*{1.2cm}

{\color{titleColor}\Huge\bfseries
REEM V2\\[0.3cm]
Inventory Intelligence
\par}

\vspace{0.5cm}

{\color{blueColor}\Large
Pharmacy Supply Chain Intelligence
\par}

\vspace{0.3cm}

{\large
Research and Development Framework
\par}

\vspace{1.2cm}

\begin{tcolorbox}[
    colframe=titleColor,
    colback=white,
    width=0.90\textwidth,
    arc=3mm
]
\centering
\large
\textbf{From Inventory State to Management Intelligence}
\end{tcolorbox}

\vfill

{\large
September 2026
\par}

\vspace{0.4cm}

{\small
Synthetic Development Environment\\
Operational Use Prohibited
\par}

\end{titlepage}

% ============================================================
\section{Executive Summary}
% ============================================================

REEM V2 is designed as an intelligence layer above pharmacy
operational systems. Its purpose is not simply to report stock
balances, but to transform operational evidence into analytical
information that can support inventory decisions.

The inventory intelligence layer integrates four major evidence
streams:

\begin{enumerate}[leftmargin=*]
    \item sales and demand history,
    \item procurement history,
    \item delivery and lead-time history,
    \item current inventory state.
\end{enumerate}

The current development implementation successfully integrates
forecast demand with observed delivery lead times and current
inventory-batch quantities.

Fifty medicines were analysed.

The current synthetic environment produced:

\begin{itemize}
    \item 50 medicines analysed,
    \item 46 medicines with observed delivery lead-time evidence,
    \item 4 medicines using an explicit seven-day development assumption,
    \item no negative current-stock values,
    \item positive reorder points for all medicines,
    \item 50 medicines classified as ADEQUATE under the current
          development decision rule.
\end{itemize}

The result should not be interpreted as evidence that a real pharmacy
has adequate stock. The dataset is explicitly synthetic and exists
only for research and development.

% ============================================================
\section{Purpose}
% ============================================================

The purpose of the inventory intelligence model is to answer a
management question:

\begin{tcolorbox}[
    colframe=blueColor,
    colback=white,
    arc=2mm
]
\centering
\textbf{How much inventory is available relative to the amount
required to cover expected demand and replenishment uncertainty?}
\end{tcolorbox}

This requires the system to connect:

\[
\text{Demand}
\rightarrow
\text{Lead Time}
\rightarrow
\text{Replenishment Requirement}
\rightarrow
\text{Current Stock}
\rightarrow
\text{Decision}.
\]

The model therefore represents a transition from operational
record-keeping toward evidence-based inventory intelligence.

% ============================================================
\section{Analytical Architecture}
% ============================================================

The REEM V2 inventory layer is positioned between the operational
data environment and management decision-making.

\begin{figure}[H]
\centering

\begin{tikzpicture}[
node distance=0.8cm and 0.8cm,
box/.style={
    rectangle,
    rounded corners,
    draw=blueColor,
    minimum width=3.0cm,
    minimum height=0.8cm,
    align=center
},
decision/.style={
    rectangle,
    rounded corners,
    draw=titleColor,
    minimum width=3.2cm,
    minimum height=0.9cm,
    align=center,
    fill=white
},
arrow/.style={
    -{Latex[length=3mm]},
    thick
}
]

\node[box] (sales) {Sales\\Demand History};
\node[box, right=of sales] (delivery) {Deliveries\\Lead Time};
\node[box, right=of delivery] (stock) {Inventory Batches\\Current State};

\node[box, below=of sales, xshift=3.2cm] (forecast)
{Demand Forecast};

\node[box, below=of forecast] (rop)
{Reorder Point\\and Safety Stock};

\node[decision, below=of rop]
{Inventory Intelligence\\Decision};

\draw[arrow] (sales) -- (forecast);
\draw[arrow] (delivery) -- (rop);
\draw[arrow] (stock) |- (rop);
\draw[arrow] (forecast) -- (rop);
\draw[arrow] (rop) -- (decision);

\end{tikzpicture}

\caption{REEM V2 inventory intelligence architecture.}
\end{figure}

% ============================================================
\section{Evidence Architecture}
% ============================================================

A major methodological improvement in REEM V2 is the distinction
between transactional evidence and state evidence.

\begin{figure}[H]
\centering

\begin{tikzpicture}[
node distance=0.7cm and 1.0cm,
box/.style={
    rectangle,
    rounded corners,
    draw=blueColor,
    minimum width=3.3cm,
    minimum height=0.75cm,
    align=center
},
state/.style={
    rectangle,
    rounded corners,
    draw=purpleColor,
    minimum width=3.3cm,
    minimum height=0.75cm,
    align=center
},
arrow/.style={
    -{Latex[length=3mm]},
    thick
}
]

\node[box] (sales) {Sales};
\node[box, below=of sales] (purchases) {Purchases};
\node[box, below=of purchases] (deliveries) {Deliveries};
\node[box, below=of deliveries] (movements) {Stock Movements};

\node[state, right=of purchases] (batches)
{Inventory Batches\\Remaining Quantity};

\node[box, right=of batches] (evidence)
{Current Inventory\\State};

\draw[arrow] (sales) -- ++(1.2,0) |- (evidence);
\draw[arrow] (purchases) -- (batches);
\draw[arrow] (deliveries) -- ++(1.2,0) |- (evidence);
\draw[arrow] (movements) -- ++(1.2,0) |- (evidence);
\draw[arrow] (batches) -- (evidence);

\end{tikzpicture}

\caption{Separation of transactional and state evidence.}
\end{figure}

This distinction is important.

Stock movements describe inventory events such as sales and receipts.
They do not necessarily constitute a complete opening-to-current
inventory ledger.

For the present synthetic environment, current stock is therefore
derived from:

\[
\text{Current Stock}_{i}
=
\sum_b
\text{RemainingQuantity}_{ib}.
\]

This avoids incorrectly producing negative inventory merely because
the movement table does not contain a complete opening balance.

% ============================================================
\section{Demand Input}
% ============================================================

The inventory model receives expected daily demand from the REEM V2
demand forecasting layer.

For medicine $i$:

\[
\widehat{D}_i
=
\text{Selected Forecast Units per Day}.
\]

The current demand model evaluates four development forecasting
methods:

\begin{enumerate}
    \item Naive,
    \item seven-day moving average,
    \item seven-day weighted moving average,
    \item linear trend.
\end{enumerate}

The selected model is based on chronological validation performance.

The inventory model does not independently reinterpret the demand
forecast. It consumes the selected forecast as an analytical input.

% ============================================================
\section{Lead-Time Evidence}
% ============================================================

Lead time is derived from delivery records whenever delivery history
is available.

For medicine $i$:

\[
L_i =
\text{observed delivery lead time}.
\]

The current development dataset contains:

\[
46/50 = 92\%
\]

of medicines with observed lead-time evidence.

Four medicines lack sufficient delivery evidence:

\begin{itemize}
    \item D0015 -- Bisoprolol,
    \item D0023 -- Clopidogrel,
    \item D0004 -- Ambroxol,
    \item D0036 -- Insulin.
\end{itemize}

For these medicines, the development model uses:

\[
L_i = 7\text{ days}
\]

with the explicit status:

\[
\texttt{ASSUMED\_7\_DAYS}.
\]

The assumption is deliberately retained in the output rather than
being presented as observed evidence.

% ============================================================
\section{Demand During Lead Time}
% ============================================================

Expected demand during replenishment lead time is calculated as:

\[
D_{LT,i}
=
\widehat{D}_i L_i.
\]

For example, D0027 (Enalapril) has:

\[
\widehat{D}=7.8571
\]

units per day and:

\[
L=9.52
\]

days.

Therefore:

\[
D_{LT}
=
7.8571\times9.52
\approx74.8.
\]

This represents expected demand during the replenishment period.

% ============================================================
\section{Safety Stock}
% ============================================================

The current development implementation uses validation forecast
error as a provisional proxy for demand uncertainty.

The development formulation is:

\[
SS_i
=
1.65
\times
MAE_i
\times
\sqrt{L_i}.
\]

The factor 1.65 corresponds approximately to a one-sided 95\%
normal-theory service factor.

However, this formulation is explicitly provisional.

A production-quality model should eventually estimate uncertainty
directly from longitudinal demand residuals and incorporate
lead-time variability.

Therefore the current safety-stock calculation should be regarded as:

\begin{tcolorbox}[
    colframe=goldColor,
    colback=white,
    arc=2mm
]
\centering
\textbf{Development-level safety-stock estimation}
\end{tcolorbox}

rather than a final pharmaceutical inventory policy.

% ============================================================
\section{Reorder Point}
% ============================================================

The current reorder point is:

\[
ROP_i
=
D_{LT,i}+SS_i.
\]

For D0027:

\[
D_{LT}=74.8
\]

and:

\[
SS=9.87.
\]

Therefore:

\[
ROP
=
74.8+9.87
=
84.67.
\]

The current stock is:

\[
11\,007.
\]

Thus the available stock is substantially above the calculated
reorder threshold.

% ============================================================
\section{Inventory Decision Logic}
% ============================================================

The current development decision rule is:

\[
\text{Decision}_i =
\begin{cases}
\text{STOCKOUT\_RISK}, & Stock_i \leq 0\\
\text{REORDER}, & Stock_i \leq ROP_i\\
\text{WATCH}, & Stock_i \leq 1.5ROP_i\\
\text{ADEQUATE}, & Stock_i > 1.5ROP_i.
\end{cases}
\]

This is a transparent development rule.

It is intentionally simple and interpretable.

It should not yet be considered a validated operational pharmacy
policy.

% ============================================================
\section{Current Inventory Results}
% ============================================================

The validated development run produced:

\begin{center}
\begin{tabular}{lr}
\toprule
\textbf{Measure} & \textbf{Result}\\
\midrule
Medicines analysed & 50\\
Minimum current stock & 175\\
Maximum current stock & 11,619\\
Minimum reorder point & 0.77\\
Maximum reorder point & 84.67\\
Observed lead-time medicines & 46\\
Assumed lead-time medicines & 4\\
Stockout risk & 0\\
Reorder & 0\\
Watch & 0\\
Adequate & 50\\
\bottomrule
\end{tabular}
\end{center}

All medicines therefore satisfy:

\[
Stock_i > ROP_i
\]

in the current synthetic environment.

% ============================================================
\section{Stock-to-Reorder-Point Analysis}
% ============================================================

The ratio:

\[
R_i=
\frac{Stock_i}{ROP_i}
\]

provides a simple measure of how far current stock lies above the
calculated reorder threshold.

The lowest ten ratios observed were:

\begin{center}
\begin{tabular}{lrr}
\toprule
\textbf{Medicine} & \textbf{Stock} & \textbf{Stock/ROP}\\
\midrule
Cefixime & 470 & 40.1\\
Omeprazole & 526 & 43.4\\
Bisoprolol & 492 & 70.4\\
Doxycycline & 686 & 70.6\\
Spironolactone & 335 & 85.7\\
Furosemide & 543 & 90.2\\
Azithromycin & 953 & 94.8\\
Ciprofloxacin & 805 & 100.0\\
Amoxicillin & 962 & 100.1\\
Captopril & 1138 & 112.3\\
\bottomrule
\end{tabular}
\end{center}

Even the lowest observed ratio is approximately 40 times the
reorder point.

This explains why no medicine is classified as REORDER or WATCH.

% ============================================================
\section{Interpretation}
% ============================================================

The result demonstrates that the inventory intelligence pipeline is
functionally coherent.

The system successfully connects:

\[
\text{Forecast}
\rightarrow
\text{Lead Time}
\rightarrow
\text{Safety Stock}
\rightarrow
\text{ROP}
\rightarrow
\text{Current Stock}
\rightarrow
\text{Decision}.
\]

The result also demonstrates an important analytical principle:

\begin{tcolorbox}[
    colframe=purpleColor,
    colback=white,
    arc=2mm
]
\centering
\textbf{The decision engine should describe the data, not manufacture
a desired management result.}
\end{tcolorbox}

The synthetic pharmacy contains substantially more stock than is
required by the current demand and replenishment assumptions.

Consequently, the correct result is not to artificially create
reorder alerts, but to report the observed state:

\[
\boxed{\text{50 medicines: ADEQUATE}}
\]

within the current development environment.

% ============================================================
\section{Evidence Quality}
% ============================================================

REEM V2 explicitly distinguishes evidence quality.

\begin{center}
\begin{tabular}{lll}
\toprule
\textbf{Component} & \textbf{Evidence} & \textbf{Status}\\
\midrule
Demand forecast & Historical sales & Development\\
Lead time & Delivery history & 92\% observed\\
Lead time & Missing history & 8\% assumed\\
Current stock & Batch remaining quantity & Development\\
ROP & Calculated & Development\\
Decision & Rule-based & Development\\
\bottomrule
\end{tabular}
\end{center}

The evidence ladder is:

\[
\text{Calculated}
\rightarrow
\text{Analytically Useful}
\rightarrow
\text{Validated}
\rightarrow
\text{Operationally Trusted}.
\]

The current inventory model is between the first two stages and is
being developed toward formal validation.

% ============================================================
\section{Model Integrity Validation}
% ============================================================

The current validation produced the following results:

\begin{center}
\begin{tabular}{ll}
\toprule
\textbf{Integrity Test} & \textbf{Result}\\
\midrule
50 medicines & PASS\\
No negative stock & PASS\\
Positive reorder points & PASS\\
All lead times observed & FAIL\\
No missing decisions & PASS\\
Development boundary preserved & PASS\\
Operational use prohibited & PASS\\
\bottomrule
\end{tabular}
\end{center}

The single failed test is intentional evidence rather than a hidden
error.

Four medicines do not have observed delivery histories. The model
therefore correctly identifies their lead time as assumed rather
than falsely labeling it as observed.

% ============================================================
\section{Important Limitations}
% ============================================================

The current model has several important limitations.

\subsection{Synthetic Data}

All current records originate from:

\[
\texttt{SIMULATED\_REEM\_V2}.
\]

They are development records and are not operational evidence.

\subsection{Safety Stock}

The present safety-stock formulation uses forecast MAE as a proxy
for uncertainty.

A future version should use empirical demand variability and
lead-time variability.

\subsection{Inventory State}

Current stock is derived from batch remaining quantities. This is
appropriate for the present synthetic state representation but does
not replace a fully reconciled production inventory ledger.

\subsection{Decision Thresholds}

The current decision thresholds are development rules. They require
validation against actual pharmacy service levels, stockout costs,
criticality, procurement constraints, and management policy.

\subsection{Overstocked Synthetic Environment}

The current synthetic inventory contains substantially more units
than required by modeled demand.

Therefore the present results cannot be used to estimate realistic
pharmacy stock requirements.

% ============================================================
\section{Next Development Stage}
% ============================================================

The next stage should improve the scientific quality of the inventory
model rather than merely add more output.

The development roadmap is:

\begin{figure}[H]
\centering

\begin{tikzpicture}[
node distance=0.65cm,
box/.style={
    rectangle,
    rounded corners,
    draw=blueColor,
    minimum width=10cm,
    minimum height=0.72cm,
    align=center
},
arrow/.style={
    -{Latex[length=3mm]},
    thick
}
]

\node[box] (a) {Validated Current Stock State};
\node[box, below=of a] (b) {Empirical Demand Variability};
\node[box, below=of b] (c) {Lead-Time Variability};
\node[box, below=of c] (d) {Improved Safety Stock};
\node[box, below=of d] (e) {Service-Level Calibration};
\node[box, below=of e] (f) {Stockout Risk Modelling};
\node[box, below=of f] (g) {Inventory Optimisation};

\draw[arrow] (a) -- (b);
\draw[arrow] (b) -- (c);
\draw[arrow] (c) -- (d);
\draw[arrow] (d) -- (e);
\draw[arrow] (e) -- (f);
\draw[arrow] (f) -- (g);

\end{tikzpicture}

\caption{Proposed inventory intelligence development pathway.}
\end{figure}

Future extensions may include:

\begin{itemize}
    \item empirical demand standard deviation,
    \item lead-time standard deviation,
    \item service-level targets,
    \item medicine criticality,
    \item stockout probability,
    \item supplier reliability,
    \item expiry risk,
    \item ABC classification,
    \item EOQ,
    \item Monte Carlo inventory simulation,
    \item multi-objective inventory optimisation.
\end{itemize}

% ============================================================
\section{Management Intelligence Perspective}
% ============================================================

REEM V2 should ultimately move beyond the question:

\begin{quote}
``How many units are currently in stock?''
\end{quote}

toward questions such as:

\begin{itemize}
    \item What demand is expected?
    \item How long will replenishment take?
    \item How much inventory is required during replenishment?
    \item How uncertain is the requirement?
    \item Which medicines are approaching risk?
    \item Which medicines are excessively stocked?
    \item Which suppliers create replenishment risk?
    \item Which inventory should receive management attention first?
\end{itemize}

This represents the central transition from an operational pharmacy
system to a management intelligence system.

% ============================================================
\section{Conclusion}
% ============================================================

The REEM V2 inventory intelligence layer has successfully established
a coherent development pipeline linking demand forecasting,
procurement lead time, safety stock, reorder point, current inventory
state, and inventory decision.

The current validation confirms:

\begin{itemize}
    \item current stock is non-negative,
    \item reorder points are positive,
    \item delivery lead times are explicitly distinguished as observed
          or assumed,
    \item all 50 medicines receive a decision,
    \item synthetic-data boundaries are preserved,
    \item no operational conclusion is claimed.
\end{itemize}

The model is therefore ready for the next stage of research:
improving uncertainty estimation, service-level modelling, stockout
risk prediction, supplier intelligence, and inventory optimisation.

\begin{tcolorbox}[
    colframe=titleColor,
    colback=white,
    arc=3mm
]
\centering
\large
\textbf{REEM V2 Principle}\\[0.2cm]
\textit{
Operational data should become evidence,\\
evidence should become intelligence,\\
and intelligence should support better decisions.
}
\end{tcolorbox}

\vfill

\begin{center}
\small
\textbf{REEM V2}\\
Research and Development Environment\\
Data Source: SIMULATED\_REEM\_V2\\
Data Status: DEVELOPMENT\\
Operational Use: PROHIBITED
\end{center}

\end{document}
